\par\medskip\noindent
\SGAInlineRefTarget{sga3:xix:proposition:3-8}\textbf{\textit{\SGARefLink{sga3:xix:proposition:3:8}{Proposition 3.8}.}}---
\textit{Let \(S\) be a scheme, let \(G\) be an \(S\)-group, and let
\(T\) be a maximal torus of \(G\).  The functor \(\mathcal R\) of
roots of \(G\) with respect to \(T\) is represented by a finite
twisted constant \(S\)-scheme \textup{(\SGARefLink{sga3:x:definition:5-1}{Exposé X, 5.1})} which is an
open and closed subscheme of}
\[
  \underline{\operatorname{Hom}}_{S\text{-gr.}}
  (T,\mathbb G_{m,S}).
\]

\textit{A subset
\(R\subset
\operatorname{Hom}_{S\text{-gr.}}(T,\mathbb G_{m,S})\)
is a root system of \(G\) with respect to \(T\) if and only if the
morphism}
\[
  R_S
  \longrightarrow
  \underline{\operatorname{Hom}}_{S\text{-gr.}}
  (T,\mathbb G_{m,S}),
\]
\textit{defined by the preceding inclusion, induces an isomorphism}
\[
  R_S\xrightarrow{\ \sim\ }\mathcal R.
\]
\SGARefTarget{sga3:xix:proposition:3:8}

\par\medskip\noindent
\textbf{\SGAInlineRefTarget{sga3:xix:paragraph:3-9}3.9.}---
Let \(S\) be a scheme, let \(G\) be a reductive \(S\)-group, let \(T\)
be a maximal torus of \(G\), and let \(\alpha\) be a root of \(G\)
with respect to \(T\), that is, a section of \(\mathcal R\).  Consider
the kernel \(\operatorname{Ker}(\alpha)\) of \(\alpha\), its unique
maximal torus \(T_\alpha\), and the centralizer of the latter,
\[
  Z_\alpha=\underline{\operatorname{Centr}}_G(T_\alpha).
\]
This is a closed \(S\)-subgroup of \(G\), reductive (\SGARefLink{sga3:xix:paragraph:2:8}{2.8}), and of
semisimple rank one (\SGARefLink{sga3:xix:proposition:1-12}{1.12}).  Moreover,
\[
  \operatorname{Lie}(Z_\alpha/S)
  =
  \mathfrak t
  \oplus\mathfrak g^\alpha
  \oplus\mathfrak g^{-\alpha};
\]
hence \(\{\alpha,-\alpha\}\) is a root system of \(Z_\alpha\) with
respect to \(T\).
\SGARefTarget{sga3:xix:paragraph:3:9}

\par\medskip\noindent
\textbf{3.10.}---
Let \(S\) be a scheme, let \(G\) be a reductive \(S\)-group, let \(T\)
be a maximal torus of \(G\), and let \(\alpha\) be a root of \(G\)
with respect to \(T\).  If \(q\in\mathbb Q\) and \(q\alpha\) is a root
of \(G\) with respect to \(T\), then \(q=1\) or \(q=-1\).  This follows
at once from \SGARefLink{sga3:xix:proposition:1-12}{1.12}.
\SGARefTarget{sga3:xix:paragraph:3:10}
